\scanpage{35}
est donné par
\[
 S_i\simeq I\setminus\{i\}\tag*{(8)}
\]
\margin{\unsure{ens|un} des \inserted{deux} sommets sur la droite}

L'ens des deux orientations de $D_i$ considéré comme support d'un bigône) est donc
\[
 \Omega_i\simeq S_i\wedge A_i\simeq
 \underbrace{(I\setminus\{i\})\wedge
 \underbrace{\bigwedge_{\alpha\in I\setminus\{i\}}\Phi_\alpha}_{\Phi_i\wedge\operatorname{or}(I)}}
 _{\substack{\Phi_i\wedge\bigwedge_{\alpha\in I\setminus\{i\}}(I\setminus\{\alpha\})\\\text{canoniquement}}}\tag*{(9)}
\]
\margin{à cause de (6bis)}

Notons que
\[
\begin{aligned}
 \underbrace{\bigwedge_{i\in I} A_i}_{\shortparallel}&\simeq\mathbb{1}\\
 \bigwedge_{\substack{\alpha,i\in I\\\alpha\ne i}}\Phi_\alpha
 &=\bigwedge_{\alpha\in I}\Phi_\alpha^2\simeq\mathbb{1}
\end{aligned}\tag*{(10)}
\]
d'où
\[
 \bigwedge_{i\in I}\Omega_i\overset{\varphi}{\simeq}
 \bigwedge_{i\in I}S_i
 \left[\simeq\bigwedge_{i\in I}(I\setminus\{i\})
 \underset{\substack{\text{isom can.}\\\text{évident associé}\\\text{à l'ens. à 3}\\\text{éléments }I}}{\overset{\sim}{\leftarrow}}
 \operatorname{or}(I)\right]\tag*{(11)}
\]

\uline{NB} \deleted{Quand} on part d'une famille \deleted{\unsure{de|i.e.}} de trois bigônes $(S_i,A_i)_{i\in I}$, la donnée d'un isom (10) équivaut à celle d'un isom $\varphi\unsure{\}|)}$ (11) (qui, pour $S_i=I\setminus\{i\}\quad\forall i\in I\}$, se réduit à un isom
\[
 \mathop{\deleted{\unsure{\pi}}\bigwedge}\limits_{i\in I}^{\unsure{i'}}\Omega_i\simeq\operatorname{or}(I)\tag*{(11 bis)}
\]
ou encore à un isom
\[
 \left(\bigwedge_{i\in I}^{\unsure{\epsilon|{\,}'}}\Omega_i\right)\wedge\operatorname{or}(I)\simeq\mathbb{1}\tag*{(12)}
\]
% archive sheet number: 34
